% ---------------------------------------------------------------------------
% Suggested replacement text for App. F (forward step, inner step, truncation,
% Algorithm 1) so that the appendix matches the code (cogent/solver.py).
% Equation numbers (104)-(123) refer to the current draft; adapt the macros.
% Two versions of the F.3 implementation paragraph are given: (A) the default
% block-coordinate step (enhancement="bcd", used for the main figure) and (B) the
% Frank-Wolfe step already written in F.3 (enhancement="fw"), with the two
% details the code needs to run it.  Use one of them.
% ---------------------------------------------------------------------------

% ===== F.3, first paragraph (replace "With the active set fixed, Phi_k is convex ...") =====
With the active set fixed we decrease $\Phi_k$ by block-coordinate descent over the rows of
$\Omega$ and over $\{\theta_i\}$. Note that $\Phi_k$ is convex in $\{\theta_i\}$ for fixed $\Omega$,
but not in $\Omega$: for a non-edge ($\bar A_{uw}=-1$) the graph term contributes the concave
function $\sqrt{\omega_{u,i}\,\omega_{w,i}}$. Convexity of (14) holds in the atomic weights over the
fixed (infinite) atomic set, not in the reduced variables. As in the enhancement step of CoGEnT,
any update that keeps every row of $\Omega$ in $\Delta^k$ and does not increase $\Phi_k$ can be used.

% ----- (A) implementation of the inner step, enhancement="bcd"
%       (replace the last sentence of the F.3 intro) -----
In our implementation a row $\Omega_{v,:}$ is replaced by the best of the following candidates:
the Frank--Wolfe vertex $\delta_{i^\star_v}$ of (109) and the other vertices with the smallest
partial derivatives (105), the Frank--Wolfe point $(1-\alpha)\Omega_{v,:}+\alpha\,\delta_{i^\star_v}$,
with $\alpha\in\{\tfrac14,\tfrac12\}$, and the pairwise steps that move half or all of the mass of
the active atom with the largest partial derivative to $i^\star_v$. Only row and column $v$ of $L$ and the parameter $\theta_v$ change, so the
change of $\Phi_k$ is evaluated exactly in $O(n(k+d))$ operations. The row is updated only if
$\Phi_k$ decreases. The atoms are updated by a majorize--minimize step for the fusion term,
$\|z\|_2\le \|z\|_2^2/(2\|z_0\|_2)+\|z_0\|_2/2$. For the loss (22) this step is a $k\times k$
linear system for the feature block and a simplex-constrained quadratic program for the
label block; for a general $f_v$ it is solved by accelerated projected gradient.

% ----- (B) alternative: keep the Frank--Wolfe enhancement step of F.3, enhancement="fw"
%       (add after (113)) -----
The linear minimization (112) over $\Theta_f$ requires $\Theta_f$ to be compact. For the loss (22)
we use $\Theta_f=\{a:\|a\|_2\le R\}\times\Delta^K$ with $R=\max_v\|z_v\|_2$. This does not change
the optimal value of (14): projecting every feature block $a_v$ onto the convex hull of the
features does not increase the loss or the fusion term. The minimizer in (112) is then
$a=-R\,g_a/\|g_a\|_2$ and the vertex of $\Delta^K$ with the smallest entry of $g_\pi$. In (105)
the graph part $-(\bar A e_i)_v/e_{i,v}$ is unbounded as $e_{i,v}\to0$; we replace $e_{i,v}$ by
$\max(e_{i,v},\varepsilon)$. Both updates (111) and (113) use one common step size per block,
chosen by a line search on $\Phi_k$, and are applied only if $\Phi_k$ decreases. Because the graph
part of (105) is of order $1/\varepsilon$ on every empty entry, a small $\varepsilon$ lets the
graph term dominate the row direction even when $\mu$ is large; the exact-change row updates of
(A) do not have this sensitivity.

% ===== F.4, replace the paragraphs "Solving the oracle", "Forward update" and
%       "Implementation used in Sec. 6" =====
\paragraph{Solving the oracle.}
Along (114) node $v$ can give at most its whole mass, so the admissible step
$\rho\le 1/\max_v e_v^2$ depends on the candidate $e$. As in the conditional-gradient step of
CoGEnT, where every atom is compared at the boundary of the feasible set, we compare candidates
by the first-order change of $\Phi$ at their maximal step, i.e.\ by the rate (115) multiplied
by $1/\max_v e_v^2$. Writing $u_v := e_v^2/\max_w e_w^2\in[0,1]$ for the fraction of the mass
of node $v$ that is moved to the new atom, this quantity reads
\begin{equation}
  \Delta(u,\theta) \;=\; -\sum_{v,w\in\mathcal V}\bar A_{vw}\sqrt{u_v u_w}
  \;+\;\sum_{v\in\mathcal V} u_v\, d_v(\theta),
  \qquad u\in[0,1]^n,\ \theta\in\Theta_f .
  \label{eq:oracle}
\end{equation}
We search for candidates with $u=\mathbf 1_S$, i.e.\ every node of a set $S$ moves the same
fraction of its mass. For fixed $\theta$, \eqref{eq:oracle} is then a quadratic function of
$\mathbf 1_S$, and for fixed $S$ the best $\theta$ is given by (121). We minimize it approximately.

Every node first chooses its own parameter $\theta^{\mathrm{new}}_v$ by (117), which gives the
scores $s_v$ of (118) and the lower bound
$\Delta(u,\theta)\ge -\sum_{v,w}\bar A_{vw}\sqrt{u_vu_w}+\sum_v u_v s_v$. With $u_v=e_v^2$, the
minimization of this bound over unit vectors is the Rayleigh quotient of $\bar A-\mathrm{diag}(s)$.
We take the nonnegative part of its leading eigenvectors (for both signs), and among their level
sets $S_\tau=\{v: e_v\ge\tau\}$ we keep the one that minimizes the bound with $u=\mathbf 1_{S_\tau}$.
Without a graph term these eigenvectors carry no information, so we also start from the single
nodes with the smallest $s_v$.

From every start we alternate two steps:
(i) $\theta^{\mathrm{new}}$ is the minimizer (121) with weights $u_v=\mathbf 1\{v\in S\}$;
(ii) $S$ is improved by single-node insertions and removals that decrease \eqref{eq:oracle} with
$d_v=d_v(\theta^{\mathrm{new}})$.
Without the graph term the second step is exact: $u_v=\mathbf 1\{d_v(\theta^{\mathrm{new}})<0\}$,
i.e.\ every node that the new atom fits better moves to it.

\paragraph{Remark (tied node parameters).}
The rate (115)--(116) uses $g_{vw}=0$ when $\theta_v=\theta_w$. Along (114) such a pair satisfies
$\theta_v(t)-\theta_w(t)=t\,(u_v-u_w)(\theta-\theta_v)$. The exact one-sided derivative therefore
contains the additional nonnegative term
$\gamma\mu\sum_{v<w:\,\theta_v=\theta_w}|u_v-u_w|\,\|\theta-\theta_v\|_2$. At hard assignments
every class is a tied group, and (115) can then even have the wrong sign. Since the rate is in any
case a first-order model, valid only for small steps, we do not rely on the single best candidate.
The several best candidates are passed to the line search below, which uses the exact $\Phi$. Using
the exact (tie-aware) rate for the ranking instead was not better in our experiments.

\paragraph{Forward update.}
The new atom $(u,\theta^{\mathrm{new}})$ is appended along (114) with $\rho\, e_v^2 = t\,u_v$:
\begin{equation}
  \Omega^{(t+\frac12)}_{v,i}=
  \begin{cases}(1-t\,u_v)\,\omega^{(t)}_{v,i}, & i\le k,\\ t\,u_v, & i=k+1,\end{cases}
  \qquad \theta^{(t+\frac12)}_{k+1}=\theta^{\mathrm{new}},\qquad k\leftarrow k+1,
\end{equation}
where, among the six candidates with the most negative rate \eqref{eq:oracle}, the candidate and
$t\in[0,1]$ minimize $\Phi_{k+1}$ along this path (line search, as in CoGEnT). Every row of $\Omega$
stays in the simplex. The line search does not increase the objective. If it returns $t=0$ for
every candidate, or no candidate has a negative rate, no atom is added. The inner step of
Sec.~F.3 is then run on the enlarged active set.

% ===== F.5, candidate selection (add after the definition of Phi^{(-j)}) =====
As in Algorithm~2 of CoGEnT, a cheap screen selects the candidate: we evaluate
$\Phi^{(-j)}$ without re-optimization for the lightest atoms and for the atoms closest to
another atom, and re-optimize only the best candidate before testing (123).

% ===== Algorithm 1, lines 1-7 =====
% 1: Initialize k <- 1, Omega <- 1_n, theta_1 <- argmin_{theta in Theta_f} sum_v f_v(z_v; theta)
% 2: Phi_prev <- Phi_1(Omega, {theta_1})
% 3: for t = 1, ..., T do
%      // (a) Forward step: select a new atom with the oracle (Sec. F.4)
% 4:   Compute theta_v^new by (117), s_v by (118), and the leading eigenvectors of Abar - diag(s)
% 5:   For each start (level sets of the eigenvectors, single nodes with smallest s_v):
%      alternate theta^new by (121) and single-node flips of S on (eq:oracle)
% 6:   Line search along (114) on the six candidates with the most negative rate; append the
%      candidate with the lowest Phi_{k+1} if it decreases Phi; k <- k+1
%      // (b) Inner step ... (unchanged)
%      // (c) Backward step ... (unchanged)
%    Stop when no candidate decreases Phi and the inner/backward steps make no progress.
